\section{Introduction}
\label{sec:intro}

\IEEEPARstart{T}{emporal} data, systematically recorded as univariate or multivariate time series, serves as the fundamental perceptual substrate for understanding dynamic physical, economic, and biological systems. Across modern industries, time series forecasting, anomaly detection, imputation, and classification drive mission-critical workflows: stabilizing renewable energy power grids, optimizing multi-echelon retail supply chains, executing high-frequency financial risk arbitrage, predicting clinical sepsis in intensive care units, and monitoring aerospace telemetry~\cite{jin2023largemodels}.

Historically, time series modeling evolved through three distinct algorithmic epochs:
\begin{enumerate}[leftmargin=*]
    \item \emph{Statistical and Econometric Formulation (1970s--2010s)}: Dominated by classical stochastic processes, such as Auto-Regressive Integrated Moving Average (ARIMA), Exponential Smoothing (ETS), Vector Autoregression (VAR), and Gaussian Processes. While mathematically rigorous and interpretable, these methods rely on linear assumptions and fail to capture complex nonlinear dependencies across cross-domain series.
    \item \emph{Task-Specific Deep Neural Networks (2015--2022)}: Catalyzed by recurrent architectures (LSTMs, GRUs), temporal convolutional networks (TCNs), and specialized Transformer variants (Informer, Autoformer, FEDformer, and PatchTST~\cite{nie2023patchtst}). Although these architectures captured long-range nonlinear temporal correlations, each model was trained in isolation on a single, isolated dataset (e.g., ETT, Weather, Electricity), exhibiting virtually zero out-of-distribution generalization.
    \item \emph{The Foundation Model Epoch (2022--Present)}: Propelled by the extraordinary success of pre-trained foundation models in Natural Language Processing (e.g., GPT-4, Llama) and Computer Vision, the temporal machine learning community has initiated a paradigm shift toward universal, cross-domain foundation models~\cite{ansari2024chronos,das2024timesfm,woo2024moirai}. Rather than training narrow models from scratch for every new dataset, practitioners now seek models capable of strong zero-shot and few-shot transfer across unseen domains, frequencies, and channels.
\end{enumerate}

\subsection{Unique Challenges of Temporal Signals}
Transferring the foundation model paradigm to time series is non-trivial due to stark physical and mathematical differences between temporal signals and natural language:
\begin{itemize}[leftmargin=*]
    \item \emph{Absence of a Discrete Natural Vocabulary}: Unlike natural text, which is composed of finite, discrete lexical tokens with stable semantic meanings, continuous time series $\mathbf{x} \in \mathbb{R}^T$ consist of real-valued scalar measurements bounded only by sensor physics.
    \item \emph{Extreme Heterogeneity in Sampling Rates and Dynamics}: Real-world time series exhibit widely divergent temporal scales—ranging from milliseconds in financial order books, to hourly power consumption, to daily retail demand, to quarterly economic indicators. A model must reconcile varying seasonalities without destructive aliasing.
    \item \emph{Variable Dimensionality and Missing Values}: Time series data dynamically switch between univariate streams and arbitrary multivariate configurations ($C \ge 1$), frequently accompanied by irregular observation intervals, sensor failure dropouts, and structural trend shifts.
\end{itemize}

\begin{table*}[t]
\caption{Comparison of This Survey with Related Surveys on Large Temporal Models.}
\label{tab:survey_comparison}
\centering
\small
\begin{tabularx}{\textwidth}{lcccccX}
\toprule
\textbf{Survey Reference} & \textbf{Year} & \textbf{PRISMA} & \textbf{Native TSFMs} & \textbf{Repurposed LLMs} & \textbf{Critiques \& Audits} & \textbf{Primary Focus \& Distinguishing Scope} \\
\midrule
Wen et al.~\cite{jin2023largemodels} & 2023 & \xmark & \xmark & \xmark & \xmark & Survey of specialized Transformer architectures on single datasets. \\
Jin et al.~\cite{jin2023largemodels} & 2024 & \xmark & Partial & \cmark & Partial & Broad survey covering both general time series and spatio-temporal graphs. \\
Gruver et al.~\cite{gruver2023llmtime} & 2023 & \xmark & \xmark & \cmark & \xmark & Empirical investigation of prompting LLMs as zero-shot forecasters. \\
Fons et al.~\cite{fons2024arellmsuseful} & 2024 & \xmark & \xmark & \cmark & \cmark & Empirical critique testing whether language backbones add genuine inductive value. \\
\textbf{This Survey} & \textbf{2026} & \cmark & \cmark & \cmark & \cmark & \textbf{Comprehensive PRISMA systematic review (2021--2026); unifying Native TSFMs, Repurposed LLMs, tokenization trade-offs, scaling laws, and leakage audits.} \\
\bottomrule
\end{tabularx}
\end{table*}

\subsection{The Bifurcated Landscape: Native TSFMs vs. LLM4TS}
To address these challenges, the research community has split into two competing architectural paradigms:
\begin{enumerate}[leftmargin=*]
    \item \emph{Native Time Series Foundation Models (TSFMs)}: Architectures designed and trained from scratch on massive multi-domain temporal corpora comprising tens to hundreds of billions of observations. Prominent exemplars include Amazon Chronos~\cite{ansari2024chronos,ansari2025chronos2}, Google TimesFM~\cite{das2024timesfm,das2024incontext}, Salesforce MOIRAI~\cite{woo2024moirai,woo2024moiraimoe,woo2025moirai2}, CMU MOMENT~\cite{goswami2024moment}, IBM TTM~\cite{ekambaram2024ttm}, Morgan Stanley/Mila Lag-Llama~\cite{rasul2023lagllama}, and the THUML Timer family~\cite{liu2024timer,liu2024timerxl,thuml2025sundial,thuml2026timers1}.
    \item \emph{Repurposed Large Language Models (LLM4TS)}: Approaches that leverage the pre-trained weights, attention layers, and sequence modeling capabilities of multi-billion parameter language models (such as Llama, GPT-2, or GPT-4). Techniques include cross-modal token reprogramming (e.g., Time-LLM~\cite{jin2023timellm}, GPT4TS/OFA~\cite{zhou2023onefitsall}, TEMPO~\cite{cao2023tempo}), direct numeric prompting with ASCII decimal strings (e.g., LLMTime~\cite{gruver2023llmtime}), and multimodal question-answering agents (e.g., Time-MQA~\cite{jin2025timemqa}, TimeOmni-1~\cite{jin2025timeomni}).
\end{enumerate}

\subsection{Key Contributions of This Survey}
While several initial surveys have examined aspects of deep learning for time series, the rapid surge of foundation models from 2024 to 2026 demands a rigorous, systematic, and up-to-date synthesis. As highlighted in Table~\ref{tab:survey_comparison}, this work delivers the following contributions:
\begin{enumerate}[leftmargin=*]
    \item \emph{PRISMA 2020 Systematic Review Protocol}: We execute a strictly reproducible review protocol across scholarly APIs, logging 83 initial records, screening against standardized criteria, and synthesizing 41 core benchmark and foundation models.
    \item \emph{Rigorous Unified Taxonomy}: We establish an orthogonal five-dimensional taxonomy deconstructing models along: (i) architectural backbone, (ii) temporal tokenization, (iii) prediction head and uncertainty formulation, (iv) adaptation mechanics, and (v) operational task scope.
    \item \emph{Deep Comparative Synthesis of Native TSFMs vs. LLM4TS}: We analyze the theoretical and empirical tensions between compact native models ($1\text{M}$--$700\text{M}$ parameters) and multi-billion parameter language adaptations ($7\text{B}$--$70\text{B}$ parameters), examining computational efficiency, inference latency, and inductive bias alignment.
    \item \emph{Benchmark Audit, Data Contamination, and Scaling Laws}: We survey modern evaluation suites (e.g., GIFT-Eval~\cite{salesforce2024gifteval}), expose systemic data leakage risks~\cite{gong2025rethinkingeval}, analyze empirical power-law scaling in temporal models, and assess probabilistic calibration reliability~\cite{li2026evaluatingreliability}.
    \item \emph{Open Research Agenda}: We articulate key open challenges, including irregular sampling equivariance, continuous flow matching, multimodal sensor fusion, and edge deployment.
\end{enumerate}

\subsection{Structure of the Survey}
The remainder of this survey is organized as follows: Section~\ref{sec:prelim} formalizes mathematical preliminaries and foundational objectives. Section~\ref{sec:taxonomy} presents our comprehensive taxonomy and comparative design table. Section~\ref{sec:native} deconstructs Native TSFMs. Section~\ref{sec:llm4ts} investigates Repurposed LLM methods. Section~\ref{sec:eval} evaluates benchmark suites, scaling laws, and critical debates. Section~\ref{sec:outlook} delineates open challenges, and Section~\ref{sec:conclusion} concludes the paper.
